\documentclass[12pt,a4paper]{article}
\usepackage[top=22mm,bottom=24mm,left=27mm,right=27mm]{geometry}
\usepackage{amsmath,amssymb,xcolor}
\definecolor{link}{HTML}{205973}
\usepackage[hypertexnames=false,colorlinks=true,linkcolor=link,urlcolor=link]{hyperref}
\hypersetup{pdftitle={Proposition 2. Two-coordinate descent for the sampled dual}}
\setlength{\parindent}{0pt}
\setlength{\parskip}{5pt}
\raggedbottom
\newcommand{\R}{\mathbb R}
\newcommand{\lead}[1]{\par\smallskip\noindent\textbf{#1}\enspace}
\begin{document}
{\LARGE\bfseries Proposition 2. Two-coordinate descent\\[2pt] for the sampled dual\par}
\medskip
At any feasible, suboptimal point of the sampled dual, moving mass between
just two coordinates decreases the objective. This gives a stopping
test based on the gradient. It also specifies which one-dimensional
quadratic the algorithm must minimize after it chooses a pair.

\section*{The program and the stopping test}
Let \(N\ge1\) count the sampled states and \(\mathcal A\) be a finite
nonempty action set. Index the draws by \(i=1,\ldots,N\), retaining
repetitions. A coordinate \(p=(i,a)\) identifies an action \(a\) at draw
\(i\), with multiplier \(\lambda_p=\lambda_{i,a}\). For \(0<\alpha<1\) and
\(\kappa>0\), consider the quadratic program
\begin{equation}
 \begin{aligned}
 \operatorname{minimize}_{\lambda}\quad&
 F(\lambda)=\tfrac12\lambda^\top Q\lambda+R^\top\lambda+S\\
 \text{subject to}\quad&
 \lambda_{i,a}\ge0,\qquad
 \sum_{a\in\mathcal A}\lambda_{i,a}\le\frac{\kappa}{N}\quad(1\le i\le N),\qquad
 \sum_{i,a}\lambda_{i,a}=\frac{1}{1-\alpha}.
 \end{aligned}
 \tag{9}\label{eq:qp}
\end{equation}
Here \(Q\) is symmetric positive semidefinite, \(R\) is a real vector,
and \(S\) is a real constant. In the sampled dual, \(Q\) is the Gram
matrix of the vectors
\(\Phi(x_i)-\alpha\sum_y p(x_i,y,a)\Phi(y)\), where \(\Phi\) maps states
into a real Hilbert space and \(p(x,y,a)\) is the transition probability
from \(x\) to \(y\) under action \(a\). Thus \(d^\top Qd\ge0\) for every
real vector \(d\), including when \(Q\) is singular.

Write \(h=Q\lambda+R\) for the gradient at \(\lambda\), and let \(e_p\)
be the vector whose sole nonzero entry is a one at \(p\). A transfer
\(d=e_p-e_q\), with \(p=(i,a)\ne q=(j,b)\), moves mass from \(q\) to \(p\).
At a feasible \(\lambda\) it permits a positive step precisely when
\begin{equation}
 \lambda_{j,b}>0,\qquad
 i=j\quad\text{or}\quad\sum_c\lambda_{i,c}<\kappa/N .
 \tag{L}\label{eq:legal}
\end{equation}

\textbf{Proposition 2.} If a feasible \(\lambda\) is suboptimal for
\eqref{eq:qp}, there are two distinct coordinates satisfying \eqref{eq:legal}
and \(h_p<h_q\). For some \(\eta>0\), every \(0<t\le\eta\) then gives a
feasible \(\lambda+t(e_p-e_q)\) with \(F(\lambda+t(e_p-e_q))<F(\lambda)\).
Consequently, a feasible point is globally optimal if and only if
\(h_p\ge h_q\) for every legal transfer.

\newpage
\section*{Proof of the two-coordinate result}
For any direction \(d\) and real \(t\), symmetry gives the exact expansion
\begin{equation}
 F(\lambda+td)-F(\lambda)=t\,h^\top d+\tfrac12t^2d^\top Qd .
 \tag{A}\label{eq:line}
\end{equation}
Choose a feasible \(\lambda'\) with \(F(\lambda')<F(\lambda)\).
Taking \(d=\lambda'-\lambda\) and \(t=1\) in \eqref{eq:line} gives
\(h^\top(\lambda'-\lambda)<0\). We must find a pair that improves this
linear cost while respecting the row caps.

\lead{Lemma 4 (an exchange of mass).} If finite real vectors \(x,y\)
satisfy \(\sum_i x_i=0\) and \(\sum_i x_i y_i<0\), then there are
indices \(i,j\) with \(x_i>0\), \(x_j<0\), and \(y_i<y_j\).
Indeed, the positive and negative coordinates of \(x\) carry the same
nonzero total mass. Dividing the strict inequality by that mass
shows that the \(x_i\)-weighted mean of \(y_i\) on the positive
coordinates is smaller than its \((-x_j)\)-weighted mean on the
negative coordinates. Some pair has the stated strict ordering.

\lead{Account for the row caps.}
Suppose every legal transfer has \(h_p\ge h_q\). In row \(i\) let
\(m_i=\min_a h_{i,a}\). A positive coordinate in that row must have
score \(m_i\), since any higher score would allow a negative-slope
transfer to a cheapest action in the same row. This argument
also applies to a full row, whose total is unchanged by the transfer.

Let \(\ell_i=\sum_a\lambda_{i,a}\), \(\ell_i'=\sum_a\lambda'_{i,a}\),
and \(x_i=\ell_i'-\ell_i\). The row changes sum to zero. Because
\(\lambda'\ge0\) and \(\lambda\) uses only cheapest actions within each row,
\[
 \sum_i m_i x_i
 \le \sum_{i,a}h_{i,a}(\lambda'_{i,a}-\lambda_{i,a})<0 .
\]
Applying Lemma~4 to \(x_i,m_i\) yields a gaining row \(i\) and a losing
row \(j\) with \(m_i<m_j\). The gaining row has room, since
\(\ell_i<\ell_i'\le\kappa/N\). The losing row has a positive coordinate
\(q\), since \(\ell_j>\ell_j'\ge0\). Its score is \(m_j\). Taking a
cheapest action \(p\) in row \(i\) gives a legal transfer with
\(h_p=m_i<m_j=h_q\), a contradiction.

\lead{A positive step.}
For the pair just found the linear term is \(t(h_p-h_q)<0\). Restrict
\(t\) to the donor mass and, for a cross-row transfer, to the room in
the receiving row. Both limits are positive. Shrinking \(t\) further
if necessary makes the quadratic term in \eqref{eq:line} smaller than
\(t(h_q-h_p)\). Conversely, any such improving pair contradicts global
optimality. This proves the proposition and the stopping test. \hfill\(\square\)

\newpage
\section*{Choosing a pair and solving its subproblem}
For a direction with at most two nonzero entries, the mass equality
forces the nonzero entries to be \(+t,-t\) with \(t>0\).
The normalization \(\|d\|_\infty\le1\) in the direction search (20)
imposes \(t\le1\). The objective in that search is the slope
\(t(h_p-h_q)\). Whenever \(h_p<h_q\), the best normalized slope for the
pair is attained at \(t=1\). The subsequent step in \(\lambda\) can be smaller.

The following two comparisons cover every legal transfer.
\begin{itemize}
 \item \textbf{Receiver in a row with room.} Choose a positive donor with
 largest \(h_q\) over all rows and a receiver with smallest \(h_p\) over
 all rows with room. Omit this comparison when all rows are full.
 \item \textbf{Receiver in a full row.} Its donor must belong to the
 same row. For each full row, compare the largest score on its positive
 coordinates with its smallest score over all actions.
\end{itemize}
In each group the candidate minimizes \(h_p-h_q\). If a group contains
a negative slope, the extreme scores are strictly ordered, so even
with ties the selected coordinates are distinct. Use the most negative
candidate. If every candidate difference is nonnegative, zero solves
(20) and Proposition~2 certifies global optimality for (9).

\lead{The corrected restricted objective.}
Let \(\mathcal B=\{p,q\}\) be the chosen pair and \(\bar{\mathcal B}\)
its complement. Keep \(\lambda_{\bar{\mathcal B}}\) fixed. Deleting the
terms independent of the variables \(u=\lambda_{\mathcal B}\) from
\(F\) leaves
\begin{equation}
 \tfrac12 u^\top Q_{\mathcal B\mathcal B}u+
 u^\top\bigl(R_{\mathcal B}
       + Q_{\mathcal B\bar{\mathcal B}}\lambda_{\bar{\mathcal B}}\bigr).
 \tag{22$'$}\label{eq:block}
\end{equation}
Indeed, the two off-diagonal blocks in the full quadratic form are
transposes and contribute equally. Their sum cancels the factor \(1/2\)
on the cross term. Formula \eqref{eq:line} gives the same correction
directly along the transfer.

In particular, for a negative-slope pair let \(t_{\max}\) be the donor
mass, reduced to the available row capacity for a cross-row move if
that is smaller. Set \(c=(e_p-e_q)^\top Q(e_p-e_q)\).
On \(0\le t\le t_{\max}\), the change in objective is
\[
 t(h_p-h_q)+\tfrac12 c\,t^2 .
\]
When \(c>0\) its minimum is at
\(\min\{t_{\max},(h_q-h_p)/c\}\). When \(c=0\) it is at \(t_{\max}\).
Either choice gives a strict decrease. The stopping test certifies the
current point whenever the search finds no improving pair.

\newpage
\section*{A sampled problem where the printed update stalls}
Take one sampled draw of state \(s\) in a two-state model with states
\(s,t\). Action \(A\) costs \(0\) and moves to \(s\). Action \(B\) costs
\(-1\) and moves to \(t\). Set \(\alpha=1/2\), \(\kappa=2\), the quadratic
regularization parameter \(\Gamma=1\), and scalar features
\(\Phi(s)=0,\ \Phi(t)=2\). Take both the sampling and relevance distributions
uniform, so their ratio at the draw is one.

The feature differences for the two actions are \(0\) and \(-1\).
The sampled mean feature vanishes, so the linear coefficient is just
\(\Gamma\) times the cost vector and the additive constant vanishes.
In coordinate order \(A,B\), the data of (9) are therefore
\[
 Q=\begin{pmatrix}0&0\\0&1\end{pmatrix},\qquad
 R=\begin{pmatrix}0\\-1\end{pmatrix},\qquad
 \lambda=(2-u,u),\quad 0\le u\le2 .
\]
This single row is full everywhere. At \(u=1/2\), moving mass from
\(A\) to \(B\) is legal and has slope \(h_B-h_A=-1/2\).

Both coordinates are in the working set, so there are no frozen terms.
The two objectives reduce to
\[
 \begin{array}{ll}
 \text{Exact restriction of (9)} & \tfrac12u^2-u
      =\tfrac12(u-1)^2-\tfrac12,\\[3pt]
 \text{Printed subproblem (22)} & u^2-u
      =(u-\tfrac12)^2-\tfrac14 .
 \end{array}
\]
The printed update returns \(u=1/2\), its unique minimum on \([0,2]\).
The corrected update moves to \(u=1\), improving the original objective
from \(-3/8\) to \(-1/2\).

\section*{Relation to the original result}
Proposition~2 holds as stated. The row-minimum argument accounts directly
for within-row transfers at saturated rows. The pair search in \(\S4.1\)
works with the empty-group convention above, which matters when every
row is full. In \(\S4.2\), equation (22) doubles both quadratic contributions
in \eqref{eq:block} while keeping \(R_{\mathcal B}\) fixed. The example shows
that solving this printed subproblem can cause premature stopping.
Using \eqref{eq:block}, or minimizing the exact transfer quadratic, gives the
decrease guaranteed by Proposition~2.

{\small\textbf{Sources.}
\href{https://moallemi.com/ciamac/papers/kernel-adp-2012.pdf\#page=18}{Bhat, Farias, Moallemi and Zheng, Proposition~2, \(\S4\) and Appendix~D}.
\href{https://github.com/nikhilpb/kadp/blob/main/lean/ActiveSet.lean}{Lean proposition and block formula},
\href{https://github.com/nikhilpb/kadp/blob/main/lean/Exchange.lean}{exchange lemma},
\href{https://github.com/nikhilpb/kadp/blob/main/lean/PairSelection.lean}{pair search},
and \href{https://github.com/nikhilpb/kadp/blob/main/lean/Examples/Proposition2.lean}{stalling example}.}
\end{document}
